\documentclass[12pt,a4paper]{article}
\usepackage[utf8]{inputenc}
\usepackage[T1]{fontenc}
\usepackage[spanish]{babel}
\usepackage[margin=1.4cm]{geometry}
\usepackage{array,tabularx}
\usepackage{enumitem}
\usepackage{titlesec}
\usepackage{xcolor}
\usepackage{hyperref}

% Restrained, high-contrast palette shared by headings, emphasis and links.
\definecolor{HeadingColor}{HTML}{243B53}
\definecolor{InstructionColor}{HTML}{704214}
\definecolor{ConceptColor}{HTML}{205C50}
\definecolor{LinkColor}{HTML}{0047AB}
\hypersetup{colorlinks=true,allcolors=LinkColor,pdfborder={0 0 0}}
\newcommand{\guidelink}[2]{\href{#1}{\textcolor{LinkColor}{\textbf{\underline{#2}}}}}
\newcommand{\instruction}[1]{\textcolor{InstructionColor}{\textbf{#1}}}
\newcommand{\concept}[1]{\textcolor{ConceptColor}{\textbf{#1}}}

\setlength{\parindent}{0pt}
\setlength{\parskip}{3pt}
\setlength{\tabcolsep}{5pt}
\renewcommand{\arraystretch}{1.12}
\setlist{leftmargin=0.6cm,topsep=3pt,partopsep=0pt,parsep=0pt,itemsep=1pt}
\titleformat{\section}{\normalsize\bfseries\color{HeadingColor}}{}{0pt}{}
\titlespacing*{\section}{0pt}{8pt}{3pt}
\newcolumntype{Y}{>{\raggedright\arraybackslash}X}
\newcommand{\checkbox}{\fbox{\rule{0pt}{0.8ex}\rule{0.8ex}{0pt}}}

\begin{document}

\begin{center}
{\Large\bfseries\color{HeadingColor} C5--C6: Organizar la evidencia y priorizar una oportunidad}\\[3pt]
Emprendimiento e Innovación --- Grado 10
\end{center}

\section*{Objetivo y punto de partida}
\concept{C5:} Organizar información contextual y relacionar usuarios, acciones, necesidades y hallazgos.\\
\concept{C6:} Priorizar una oportunidad según relevancia, viabilidad de acción y beneficio esperado.

Ya cuentan con oportunidades de C1, observaciones de C2 y evidencia de usuarios de C3--C4. Trabajen sobre la misma situación investigada: C5 organiza lo que saben; C6 decide qué aspecto merece atención. \instruction{Todavía no propongan ni seleccionen soluciones.}

\section*{Documento acumulativo e integridad de la evidencia}
Esta guía es de consulta. Continúen el \textbf{mismo documento del proyecto en LaTeX}; si aún no existe, créenlo en \guidelink{https://www.overleaf.com/}{Overleaf}. Mantengan título e integrantes, indiquen grado y fecha de actualización, y conserven los registros de C1--C4 como fuentes. Añadan C5 y C6; después incorporarán C7--C10 al mismo archivo. Entreguen el PDF completo actualizado en cada etapa.

Usen referencias internas a secciones y códigos; no copien registros entre etapas. Distingan \concept{evidencia} (lo observado o expresado), \concept{conclusión} (lo que permite inferir), \concept{supuesto} (lo no verificado) e \concept{idea} (una propuesta por explorar). No inventen citas ni fechas. Si utilizan un caso simulado autorizado por el docente, identifíquenlo y mantengan esa condición en todo el proyecto.

La IA es opcional: solo apoya código LaTeX, formato y organización de contenido propio. \instruction{No debe fabricar evidencia, conclusiones, decisiones, ideas ni aportes de integrantes.} Los ejemplos de las guías son ficticios, no evidencia del equipo.

\section*{C5 --- Organicen y conecten lo investigado}
\begin{enumerate}
\item Retomen el perfil, mapa, recorrido o lista de hallazgos existente. Complétenlo donde haga falta; no elaboren una versión nueva de cada formato. Si aún no hay una representación integrada, elijan \textbf{una}: mapa, tabla o diagrama. Pueden usar \guidelink{https://app.diagrams.net/}{diagrams.net} e incorporarla al archivo.
\item Muestren al usuario principal, qué necesita lograr, sus acciones antes, durante y después de la dificultad, y las relaciones con lugares, momentos, personas y recursos. Sustenten la elección del usuario con un registro de C2 o C3--C4.
\item Organicen los hallazgos con códigos H1, H2, etc. Si ya existe una lista, añadan códigos y fuentes allí. Usen la siguiente estructura como orientación, \textbf{no como una segunda tabla obligatoria}.
\end{enumerate}

{\small
\begin{tabularx}{\linewidth}{|p{0.08\linewidth}|Y|Y|Y|}
\hline
\textbf{Código} & \textbf{Hallazgo: conclusión} & \textbf{Evidencia y fuente} & \textbf{Relación / pendiente} \\
\hline
H1, etc. & Qué aprendieron del usuario y el contexto. & Registro de observación o respuesta; actividad y fecha, si existe. & Vínculo con otro hallazgo; supuesto o pregunta aún abierta. \\
\hline
\end{tabularx}}

Vinculen los códigos H con la representación; no repitan sus descripciones. Señalen una pregunta pendiente importante en el mismo registro. \textbf{Ejemplo ficticio:} una fila observada es evidencia; afirmar que la espera se concentra al inicio requiere comparar registros; atribuirla a falta de personal es un supuesto si no se investigó esa causa.

\concept{Evidencia de C5:} una representación contextual y hallazgos trazables a sus fuentes, integrados en los registros existentes. Estos hallazgos serán el respaldo de C6.

\newpage
\section*{C6 --- Comparen oportunidades dentro de la situación investigada}
Una \concept{oportunidad} es una necesidad o aspecto que vale la pena mejorar, no un producto. Delimiten \textbf{dos o tres aspectos de mejora} respaldados por C5. No reabran proyectos de C1 que no investigaron ni inventen problemas para completar la comparación. Si solo hay respaldo para una opción, incluyan otras como provisionales y expliciten qué falta conocer.

\textbf{Ejemplo ficticio:} disponer de tiempo para comer y comprender cuánto puede durar la espera son aspectos diferentes de una misma experiencia de compra, si los usuarios los reconocen. ``Crear una aplicación'' ya es una solución.

Identifiquen las opciones como A, B y, si corresponde, C. Describan usuario, necesidad y contexto brevemente en la matriz; relacionen cada opción con códigos H, sin rehacer los hallazgos. Comparen usando la misma escala:

{\small
\begin{tabularx}{\linewidth}{|p{0.17\linewidth}|Y|Y|Y|}
\hline
\textbf{Criterio} & \textbf{1: Bajo} & \textbf{2: Medio} & \textbf{3: Alto} \\
\hline
Relevancia & Dificultad ocasional y poco importante para los usuarios. & Recurrente o importante para algunos usuarios. & Frecuente y prioritaria para los usuarios. \\
\hline
Viabilidad de acción & Barreras importantes de tiempo, acceso o permisos. & Abordable con condiciones por confirmar. & Tiempo, acceso y recursos permiten abordarla. \\
\hline
Beneficio esperado & Efecto favorable limitado. & Mejora apreciable de la experiencia. & Mejora significativa de una necesidad importante. \\
\hline
\end{tabularx}}

\textbf{Relevancia} indica cuánto importa la situación actual; \textbf{beneficio esperado}, qué cambio favorable podría aportar atenderla. Es una expectativa razonada, no un resultado demostrado. \textbf{Viabilidad de acción} considera el alcance del grupo, acceso a usuarios, tiempo, recursos y permisos; no exige diseñar ni presupuestar una solución.

\section*{C6 --- Registren una comparación y una decisión}
Elaboren una sola matriz. Puntúen de 1 a 3; todos los criterios pesan igual. El total máximo es 9.

{\small
\begin{tabularx}{\linewidth}{|Y|c|c|c|c|}
\hline
\textbf{Opción y hallazgos H} & \textbf{Relevancia} & \textbf{Viabilidad} & \textbf{Beneficio} & \textbf{Total} \\
\hline
A: usuario, necesidad y contexto & 1--3 & 1--3 & 1--3 & 3--9 \\
\hline
B: usuario, necesidad y contexto & 1--3 & 1--3 & 1--3 & 3--9 \\
\hline
\end{tabularx}}

Debajo de la matriz, justifiquen \textbf{cada puntaje} con hallazgos H o condiciones concretas. Usen \textbf{SD (sin datos)} cuando falte respaldo: no asignen 1 por desconocimiento ni calculen un total completo para esa opción. Una opción menos investigada no es necesariamente menos relevante.

Cierren con \textbf{un párrafo de decisión}: código seleccionado, comparación con las alternativas, evidencia que respalda relevancia y beneficio, condiciones de viabilidad y limitación principal. Remitan a las justificaciones anteriores sin copiarlas. El total orienta, no sustituye el juicio: expliquen empates o elecciones con menor puntaje. Con datos insuficientes, declaren la selección provisional y el dato que falta.

\section*{Resultado y continuidad}
\concept{Evidencia de C6:} matriz justificada y una oportunidad priorizada. Conserven sus códigos y fuentes. C7 usará esta decisión para formular un desafío; no repetirá la comparación. Añadan etiquetas LaTeX a las secciones C5 y C6 para referenciarlas mediante \texttt{\textbackslash label} y \texttt{\textbackslash hyperref}.

\section*{Verificación final}
{\small\begin{itemize}[label=\checkbox]
\item C5 relaciona usuario y contexto; los hallazgos tienen fuentes, códigos y pendientes.
\item C6 compara aspectos de la misma situación y justifica los tres criterios sin ocultar datos faltantes.
\item La decisión reconoce límites y no impone una solución.
\item Añadimos C5 y C6 al mismo archivo sin duplicar registros; respetamos la integridad de la evidencia.
\item Compilamos y revisamos enlaces, tablas y legibilidad del PDF actualizado.
\end{itemize}}

\end{document}
